\documentclass[10pt]{article}
\providecommand{\FIGDIR}{fig}
\newcommand{\DIGESTDATE}{11 August 2026}
% ===================== PM Research Watch - shared preamble =====================
\usepackage[T1]{fontenc}
\usepackage[utf8]{inputenc}
\usepackage{amsmath,amssymb}
\usepackage{tgpagella}
\usepackage{tgheros}
\renewcommand{\sfdefault}{qhv}
\renewcommand{\ttdefault}{lmtt}
\usepackage{microtype}
\usepackage[english]{babel}

\usepackage[a4paper,top=17mm,bottom=17mm,left=17mm,right=17mm,
            headsep=6mm,footskip=10mm]{geometry}
\usepackage{graphicx}
\usepackage{xcolor}
\usepackage{colortbl}
\usepackage{tikz}
\usetikzlibrary{calc}
\usepackage[most]{tcolorbox}
\usepackage{titlesec}
\usepackage{fancyhdr}
\usepackage{booktabs}
\usepackage{array}
\usepackage{longtable}
\usepackage{enumitem}
\usepackage{multicol}
\usepackage{ragged2e}
\usepackage[hidelinks]{hyperref}

% ------------------------------------------------------------------ palette
\definecolor{Ink}    {HTML}{14262E}
\definecolor{Deep}   {HTML}{0F4C5C}
\definecolor{Teal}   {HTML}{1F7A8C}
\definecolor{Sky}    {HTML}{4FA3C4}
\definecolor{Sage}   {HTML}{6FA287}
\definecolor{Amber}  {HTML}{E3A008}
\definecolor{Clay}   {HTML}{B5643C}
\definecolor{Coral}  {HTML}{D1495B}
\definecolor{Violet} {HTML}{6A5B9E}
\definecolor{Slate}  {HTML}{4A6572}
\definecolor{Paper}  {HTML}{FBF9F5}
\definecolor{Mist}   {HTML}{EEEAE2}
\definecolor{Rule}   {HTML}{D9D3C9}

\pagecolor{Paper}
\color{Ink}

\hypersetup{colorlinks=true, linkcolor=Deep, urlcolor=Teal, citecolor=Teal,
            pdftitle={PM Research Watch}, pdfauthor={Automated daily digest}}

% ------------------------------------------------------------------ headings
\titleformat{\section}
  {\sffamily\bfseries\large\color{Deep}}{}{0pt}
  {}[{\vspace{1.2mm}\color{Amber}\titlerule[1.1pt]}]
\titlespacing*{\section}{0pt}{5.5mm}{2.8mm}

\titleformat{\subsection}
  {\sffamily\bfseries\normalsize\color{Teal}}{}{0pt}{}
\titlespacing*{\subsection}{0pt}{3.6mm}{1.4mm}

\setlist[itemize]{leftmargin=4.2mm, itemsep=1.1mm, topsep=1.3mm, parsep=0pt,
                  label=\textcolor{Teal}{\footnotesize$\blacktriangleright$}}

% ------------------------------------------------------------------ header/footer
\pagestyle{fancy}
\fancyhf{}
\renewcommand{\headrulewidth}{0pt}
\renewcommand{\footrulewidth}{0pt}
\fancyhead[L]{\sffamily\scriptsize\color{Slate}PM RESEARCH WATCH\;\textcolor{Rule}{\textbar}\;\DIGESTDATE}
\fancyhead[R]{\sffamily\scriptsize\color{Slate}Particulate matter monitoring \& health effects}
\fancyfoot[C]{\sffamily\scriptsize\color{Slate}\thepage}

% ------------------------------------------------------------------ components
% metric tile
\newcommand{\metric}[3]{%
  \begin{tcolorbox}[enhanced, colback=white, colframe=Rule, boxrule=0.5pt,
    arc=1.6mm, left=2.2mm, right=2.2mm, top=1.8mm, bottom=1.8mm,
    borderline west={1.7mm}{0pt}{#3}, nobeforeafter, valign=center,
    height=17mm, width=\linewidth]
    \RaggedRight
    {\sffamily\bfseries\fontsize{17}{17}\selectfont\color{#3}#1}\\[0.5mm]
    {\sffamily\fontsize{6.5}{7.4}\selectfont\color{Slate}#2}
  \end{tcolorbox}}

% pull-quote / key insight box
\newtcolorbox{keybox}[1][]{enhanced, colback=Mist, colframe=Mist,
  arc=1.6mm, left=3.2mm, right=3.2mm, top=2.4mm, bottom=2.4mm,
  borderline west={1.7mm}{0pt}{Amber}, #1}

% subtopic band
\newcommand{\band}[2]{%
  \vspace{3.2mm}\noindent
  \begin{tcolorbox}[enhanced, colback=#2, colframe=#2, arc=1.2mm,
    left=2.6mm, right=2.6mm, top=1.1mm, bottom=1.1mm, boxrule=0pt,
    width=\linewidth, nobeforeafter]
    {\sffamily\bfseries\small\color{white}#1}
  \end{tcolorbox}\vspace{1.6mm}}

% one paper entry: {short cite}{title}{journal}{doi}{body}
\newcommand{\paperentry}[5]{%
  \par\vspace{1.9mm}\noindent
  {\sffamily\bfseries\footnotesize\color{Ink}#1}\;%
  {\sffamily\fontsize{7}{8}\selectfont\color{Slate}#3}\\[0.4mm]
  {\footnotesize\itshape\color{Deep}#2}\\[0.6mm]
  {\footnotesize #4}\\[0.3mm]
  {\fontsize{6.6}{7.6}\selectfont\color{Slate}\texttt{doi:}\,\href{https://doi.org/#5}{\texttt{\detokenize{#5}}}}
  \par}

% caption style
\newcommand{\figcap}[1]{\par\vspace{1.0mm}
  {\fontsize{7.2}{8.6}\selectfont\color{Slate}\textsf{\textbf{Fig.}}\; #1\par}\vspace{1.5mm}}

\providecommand{\permil}{\textperthousand}

\setlength{\parindent}{0pt}
\setlength{\parskip}{0pt}
\renewcommand{\arraystretch}{1.18}


\begin{document}

% ==================================================================== MASTHEAD
\thispagestyle{empty}
\begin{tikzpicture}[remember picture, overlay]
  \fill[Deep] (current page.north west) rectangle
      ($(current page.north east)+(0,-67mm)$);
  \fill[Amber] ($(current page.north west)+(0,-67mm)$) rectangle
      ($(current page.north east)+(0,-68.4mm)$);
  \foreach \x/\y/\r in {22/12/0.9, 41/26/1.5, 63/9/0.7, 88/31/1.1, 112/17/1.9,
                        138/29/0.8, 159/13/1.3, 178/54/1.0, 196/21/1.6,
                        30/44/0.6, 74/57/1.2, 125/48/0.9, 168/6/0.8, 191/59/0.7,
                        52/18/1.0, 100/60/0.6, 148/20/1.1, 66/49/1.4, 205/45/0.8}
    \fill[Sky, opacity=0.30] ($(current page.north west)+(\x mm,-\y mm)$) circle (\r mm);
\end{tikzpicture}

\vspace*{4mm}
{\sffamily\fontsize{7.6}{9}\selectfont\color{Sky}%
 AUTOMATED DAILY LITERATURE DIGEST\;\textbullet\;PhD RESEARCH WATCH\par}
\vspace{2.2mm}
{\sffamily\bfseries\fontsize{27}{29}\selectfont\color{white}%
 Particulate Matter\par}
{\sffamily\fontsize{27}{29}\selectfont\color{Sky}%
 Monitoring \& Health Effects\par}
\vspace{5.0mm}
{\sffamily\fontsize{8.4}{10}\selectfont\color{white}%
 Issue 2026-08-11\;\;\textcolor{Amber}{\textbar}\;\;Entry window 11 August 2026\;\;%
 \textcolor{Amber}{\textbar}\;\;26 records carried, 24 rejected\par}
\vspace{18mm}

% -------------------------------------------------------------- metric strip
\noindent
\begin{minipage}[t]{0.190\linewidth}\metric{26}{RECORDS CARRIED\\(24 REJECTED)}{Deep}\end{minipage}\hfill
\begin{minipage}[t]{0.190\linewidth}\metric{8}{SUBTOPIC CLUSTERS\\OF TEN REPRESENTED}{Teal}\end{minipage}\hfill
\begin{minipage}[t]{0.190\linewidth}\metric{7}{RECORDS FROM CHINA\\(27\% OF CORPUS)}{Sky}\end{minipage}\hfill
\begin{minipage}[t]{0.190\linewidth}\metric{7}{EFFECT ESTIMATES\\WITH INTERVALS}{Amber}\end{minipage}\hfill
\begin{minipage}[t]{0.190\linewidth}\metric{6}{METADATA-ONLY\\RECORDS}{Clay}\end{minipage}

\vspace{6mm}

% -------------------------------------------------------------- signal of day
\section*{\sffamily\bfseries\color{Deep}Signal of the day}
\vspace{-1.5mm}{\color{Amber}\rule{\linewidth}{1.1pt}}\vspace{2.5mm}

\begin{keybox}
{\sffamily\bfseries\small\color{Deep}A national PurpleAir calibration that uses only
sensor-reported parameters --- no external meteorology, no ancillary fields --- reaches
$r=0.90$--$0.96$ against EPA reference monitors on a held-out year.}\\[1.4mm]
\footnotesize
Balasubramanian et al.\ (\textit{Environ Pollut}) train a random-forest correction for
PurpleAir PM$_{2.5}$ across the United States on 2019--2021 data and evaluate it on
\textbf{2022}, a genuine temporal hold-out rather than a random split. The design choice
that matters is the feature set: PM$_1$, PM$_{2.5}$, PM$_{10}$, temperature, pressure and
relative humidity, \emph{all reported by the sensor itself}. Nothing is drawn from a
reanalysis product or a nearby met station. Pearson $r$ against co-located EPA monitors is
\textbf{0.90--0.96} across buffer radii from 500\,m to 5\,km. Two things follow for
distributed sensing. First, a self-contained feature set is what makes a correction
deployable at the edge and in regions with no regulatory infrastructure --- the stated
motivation, and the reason this is more useful than a better-performing model that needs
a meteorological feed. Second, the buffer sweep is an implicit spatial-transferability
experiment: performance holding from 500\,m to 5\,km is the quantity a network planner
actually needs, and it is reported here as a design parameter rather than a limitation.
The caveat is that $r$ is not an error metric. No RMSE, bias or slope is quoted in the
abstract, and $r$ is insensitive to exactly the multiplicative bias that humidity-driven
hygroscopic growth produces. A correction can be strongly correlated and still wrong by a
factor. Until the error statistics are read off the paper, treat the headline as evidence
of \emph{tracking}, not of accuracy.
\end{keybox}

\vspace{3mm}

% -------------------------------------------------------------- exec brief
\section*{\sffamily\bfseries\color{Deep}What changed today}
\vspace{-1.5mm}{\color{Amber}\rule{\linewidth}{1.1pt}}\vspace{2.5mm}

\footnotesize
\begin{itemize}
\item \textbf{The forest plot is back after seven empty issues.} Seven interval estimates
were extractable today, against zero on each of 4--10 August. Both sources are
Chinese: a 24{,}333-participant thyroid cross-section and a six-year Xi'an outpatient
time series. The thyroid result --- OR \textbf{1.27 (95\% CI 1.06--1.51)} for subclinical
hypothyroidism per 10\,$\mu$g/m$^3$ PM$_{2.5}$, with FT4 down 3.57\% and TSH up 12.15\%
--- is the day's only well-powered dose--response on a novel endpoint.

\item \textbf{A negative result about lockdowns that deserves attention.} Wang et al.\
find combustion-source pollutant--respiratory associations were \emph{stronger} during
strict COVID control (SO$_2$ ER \textbf{18.40\%}, 95\% CI 10.29--27.11) than
post-pandemic (\textbf{2.91\%}, 95\% CI $-$4.82 to 11.28), while PM$_{2.5}$ and O$_3$,
which lack a dominant indoor combustion source, showed no interaction ($p>0.9$). The
authors are careful: the formal interaction tests do not reach significance
($p=0.13$, $p=0.15$) and they label it hypothesis-generating. The mechanism they
propose --- confinement shifting exposure indoors toward residential heating --- is
precisely the failure mode of ambient-monitor exposure assignment, and it is unresolvable
without personal or indoor measurement.

\item \textbf{Source-diagnostic measurement keeps beating inventories.} Chen et al.\
constrain modelled biomass-burning black carbon in southern China with year-round
$^{14}$C at three sites and recover a normalised mean bias of \textbf{$-$68\% to
$+$28\%} --- an order-of-magnitude spread in a quantity inventories report as a number.
Transboundary Southeast Asian contribution via free-tropospheric transport persists even
in winter, when continental air masses dominate and are assumed to exclude it.

\item \textbf{Emission control is redistributing nitrogen rather than removing it.}
Across China IV to VI light-duty diesel, NO falls from $>$80\% of reactive nitrogen while
NH$_3$ (SCR slip) and N$_2$O (ammonia-slip catalyst) rise. Any network measuring NO$_x$
alone is now blind to a growing fraction of what the fleet emits --- a direct argument for
speciated rather than surrogate measurement in traffic-adjacent deployments.

\item \textbf{Two exposure pathways that ambient monitoring does not see.} Indoor air and
dust supply up to \textbf{20\% of the EFSA tolerable weekly intake} of PFAS for Faroese
15-year-olds, and an urban synthetic-turf field yields fine, submicron and nano-sized
Fe/Si/Al/Br particles with Ni, Ag and Rh recovered from stub samplers worn \emph{on the
players} --- microenvironmental sampling detecting what a fixed 24-hour PM$_{2.5}$ mean of
10.93\,$\mu$g/m$^3$ reports as unremarkable.

\item \textbf{Six of 26 records ship without an abstract.} All six are Elsevier, Springer
or RSC DOIs absent from Crossref, Europe PMC and Semantic Scholar at build time. This is
the fourth consecutive issue affected and the count is rising, not falling.
\end{itemize}

\clearpage

% ==================================================================== ANALYTICS
\section*{\sffamily\bfseries\color{Deep}Corpus analytics}
\vspace{-1.5mm}{\color{Amber}\rule{\linewidth}{1.1pt}}\vspace{3mm}

\noindent\includegraphics[width=\linewidth]{\FIGDIR/f1_subtopics.png}
\figcap{Subtopic assignment of the 26 carried records; single-label by dominant
contribution. \textbf{Exposure assessment \& modelling leads at 7}, ahead of
\textbf{Sensing, forecasting \& instrumentation at 5} --- together 12 of 26, or 46\% of
the issue on the measurement side of the field. Occupational \& indoor, Burden/policy and
Respiratory \& allergic tie at 3 each; Neuro and Other clinical endpoints at 2 each;
Mechanistic toxicology at 1. \textbf{Two of the ten standing subtopics are unrepresented}
--- Cardiovascular \& metabolic and Reproductive \& developmental --- both of which
carried records as recently as last week.}

\vspace{3mm}
\noindent\includegraphics[width=\linewidth]{\FIGDIR/f2a_architecture.png}
\figcap{Study architecture across the 26 records. \textbf{Measurement campaign and
Metadata only are tied at 6 each} --- that is, as many records were carried without a
retrievable abstract as were carried on the strength of a field campaign, which is the
single most uncomfortable fact about this issue. Modelling / inventory and Review /
synthesis follow at 3 each, then Observational cross-sectional and Observational acute at
2 each, with Chamber / laboratory, Experimental / toxicology, Observational cohort and
Observational ecological at 1 apiece. \emph{Note:} the combined f2 panel is not used this
issue --- with 14 distinct endpoint labels its y-axis text overlaps the donut centre and
legend. The two panels are shown separately until that layout defect is fixed.}

\vspace{3mm}
\noindent\includegraphics[width=\linewidth]{\FIGDIR/f2b_endpoint.png}
\figcap{Health endpoint or organ system. \textbf{15 of 26 records report no health
endpoint at all} --- 8 exposure-only, 5 instrument-only, 2 emission-only --- which is the
structural signature of an instrumentation-weighted day. The remaining 11 records carry
\textbf{11 distinct endpoints, each appearing exactly once}: thyroid function, BALT and
IL-1$\alpha$/CXCL1, fronto-limbic white matter, basophil CD63, respiratory outpatient
visits, chronic cough/phlegm/COPD, haematologic malignancy, Parkinson's disease, NCD
mortality, second-hand smoke classification, and a mortality/CVD/communicable-disease
mapping review. No endpoint is replicated today, so nothing in this issue can corroborate
anything else in it.}

\clearpage

\noindent\includegraphics[width=\linewidth]{\FIGDIR/f3_metric_geography.png}
\figcap{Left --- particle metric. \textbf{The fallback bin ``PM (unspeciated) /
emissions'' is still the largest at 13 of 26}, half the corpus, despite the four
instrument-side classes added on 10 August. It is doing honest work here, because 6 of
those 13 are metadata-only records whose particle metric genuinely cannot be read, but
the bin remains the weakest column in the taxonomy. PM$_{2.5}$ only 6; composition /
speciation 3; optical properties 2; PM$_{2.5}$ and PM$_{10}$ jointly 2. Right ---
geographic provenance. \textbf{China 7 is the largest identifiable region};
``Global / multi-region'' at 8 is inflated by the six metadata-only records and the two
reviews, none of which state a site. Europe 3; Middle East \& N.\ Africa and North America
2 each; South Asia, East Asia ex-China, Sub-Saharan Africa and Oceania 1 each. Three geo
labels fell through the mapper on the first build and were pooled into the global bin;
they are now mapped (see provenance).}

\vspace{3mm}
\noindent\includegraphics[width=\linewidth]{\FIGDIR/f5_forest.png}
\figcap{The seven interval estimates in today's corpus, log scale, dashed line at the
null. \textbf{Six of seven exclude 1.0}; the exception is the post-pandemic SO$_2$
estimate, 1.029 (0.952--1.113), whose width is the whole point of that paper's phase
comparison. Metrics are heterogeneous by design and are colour-coded: OR for subclinical
hypothyroidism, GMR (geometric-mean ratio) for the two continuous thyroid biomarkers, and
RR-equiv for the four excess-risk percentages, converted as $1+\mathrm{ER}/100$ so they
sit on a common ratio axis. \textbf{They are not poolable} --- two studies, two
populations, two exposure definitions --- and no summary estimate is drawn.}

\clearpage

\noindent\includegraphics[width=\linewidth]{\FIGDIR/f6_lifecourse.png}
\figcap{Life-course windows addressed. \textbf{Working-age adulthood dominates at 5},
against 3 in puberty/adolescence, 1 in late childhood, 1 in older adults and
\textbf{zero in utero} --- the first issue since 9 August with no preconception or
prenatal record. \emph{This figure was a defect until this issue}: \texttt{plots.py}
carried a silent hardcoded fallback that rendered counts and exemplars from unrelated
earlier issues whenever the corpus JSON lacked a \texttt{LIFECOURSE} key. It is now a hard
build failure, and the counts above are read off today's records.}

\vspace{3mm}
\noindent\includegraphics[width=\linewidth]{\FIGDIR/f4_heatmap.png}
\figcap{Subtopic $\times$ study architecture; blank cells are true zeros. The two densest
cells are \textbf{Sensing $\times$ Metadata only (3)} and \textbf{Occupational \& indoor
$\times$ Measurement campaign (3)}. The first is the diagnostic one: the subtopic this
watch exists to track is also the subtopic least likely to arrive with a readable
abstract, because it publishes in Elsevier and Springer atmospheric-science titles that
deposit metadata without abstracts. Exposure assessment spreads evenly across measurement,
metadata-only and modelling (2 each) plus one review. Respiratory \& allergic is the only
health subtopic with more than one architecture at $n>1$ (2 acute, 1 cross-sectional).}

\clearpage

% ==================================================================== DIGEST
\section*{\sffamily\bfseries\color{Deep}Paper-level digest}
\vspace{-1.5mm}{\color{Amber}\rule{\linewidth}{1.1pt}}\vspace{2.5mm}

\footnotesize\noindent
Records are grouped by subtopic and ordered within group by first author. Tier A/B/C is a
judgement about how much weight the design can carry, not about the journal. Records
marked \textsc{metadata only} had no abstract in Crossref, Europe PMC, Semantic Scholar or
Scholar Gateway at build time and are listed rather than summarised.

\band{Sensing, forecasting \& instrumentation\hfill 5 records}{Amber}

\paperentry{Balasubramanian et al.\ 2026}
{Calibration model for PurpleAir-monitored PM$_{2.5}$ across the United States}
{Environ Pollut}
{Random-forest correction for PurpleAir PM$_{2.5}$ trained on 2019--2021 and evaluated on
a held-out 2022, using only sensor-reported inputs --- PM$_1$, PM$_{2.5}$, PM$_{10}$,
temperature, pressure, relative humidity --- with no external meteorology. Pearson $r$
against co-located EPA monitors is \textbf{0.90--0.96} across buffers from 500\,m to
5\,km; multiple calibrated products are released at different buffer radii. Deliberately
interpretable rather than a deep network, and explicitly aimed at regions without
regulatory infrastructure. \emph{The limitation is the metric}: $r$ measures covariation,
not agreement, and is blind to the multiplicative bias hygroscopic growth actually
produces. Without RMSE, slope or bias the accuracy claim is unverified. Still the most
directly transferable record of the day for a distributed PM network --- a national,
sensor-only, temporally validated correction is the exact artefact a deployment needs.}
{10.1016/j.envpol.2026.128935}

\paperentry{Albashir 2026}
{A physics-informed digital twin framework for health vigilance and indoor air quality
decision support}
{Res Sq (preprint)}
{Preprint proposing a physics-informed digital twin for indoor air quality in
healthcare-oriented environments, arguing that monitoring-only systems do not support
proactive decisions. \textbf{Not peer reviewed, single author, no deployment, no
validation against measured concentrations reported in the abstract.} Carried only
because the architectural claim --- coupling a physical model of a room to a sensor
stream so the system can extrapolate beyond sensor locations --- is the same idea as
physics-informed calibration on the ambient side, and the indoor literature is where it
is being tried first. Treat as a pointer to a design pattern, not as evidence. Relevance
to distributed sensing: physics-informed correction is the leading candidate for
generalising calibration away from co-location sites, and this is a low-cost read on how
that framing is developing indoors.}
{10.21203/rs.3.rs-10034990/v1}

\paperentry{\textit{Atmos Environ} 2026a \textsc{metadata only}}
{Reliability-constrained identification of primary and elevated aerosol structures from a
single tilted Mie lidar under dust-influenced conditions}
{Atmos Environ}
{\textsc{metadata only} --- no abstract in Crossref, Europe PMC or Semantic Scholar.
Listed, not summarised. Carried because the title states a specific and testable
capability: retrieving both the primary and the elevated aerosol layer from a
\emph{single tilted} Mie lidar, with an explicit reliability constraint, under dust
loading that normally defeats single-wavelength inversion. Single-instrument retrieval of
layer structure is the difference between a research lidar and one a network can afford,
and dust is the condition under which the retrieval usually fails. Whether the paper
delivers that is unknown until an abstract appears.}
{10.1016/j.atmosenv.2026.122297}

\paperentry{\textit{Atmos Environ} 2026b \textsc{metadata only}}
{Stable carbon isotopic composition of PM$_1$ emissions from biomass and coal combustion
under monitored combustion conditions}
{Atmos Environ}
{\textsc{metadata only} --- no abstract retrievable at build time. Carried because
$\delta^{13}$C signatures for PM$_1$ measured under \emph{monitored} combustion
conditions are the reference values that source-apportionment studies need and usually
borrow from decades-old literature. Biomass and coal are the two sources whose
separation matters most in residential-heating regions, and combustion-condition control
addresses the standing objection that isotopic signatures drift with burn phase. If the
reported ranges are tight this is a reference dataset; if they overlap it undermines
$\delta^{13}$C apportionment. Unresolvable without the text.}
{10.1016/j.atmosenv.2026.122296}

\paperentry{\textit{Aerosol Air Qual Res} 2026 \textsc{metadata only}}
{Identifying winter PM$_{2.5}$ sources in an urban environment using receptor modeling
and nonparametric trajectory analysis}
{Aerosol Air Qual Res}
{\textsc{metadata only} --- no abstract in Crossref, Europe PMC or Semantic Scholar.
Carried because the pairing of receptor modelling with \emph{nonparametric} trajectory
analysis is methodologically specific: it drops the functional-form assumption that
conventional trajectory statistics impose on the source--receptor relationship, which is
the assumption most likely to fail in complex winter meteorology with shallow boundary
layers. Site, sample size and factor count are all unstated. Listed for the register, not
weighted.}
{10.1007/s44408-026-00157-8}

\band{Exposure assessment \& modelling\hfill 7 records}{Amber}

\paperentry{Chen et al.\ 2026}
{Inventory uncertainty and transboundary transport of biomass-burning black carbon in
southern China}
{Environ Sci Technol}
{Year-round radiocarbon measurement of biomass-burning black carbon at two regional and
one background site in southern China, compared against tailored atmospheric-transport
simulations. The model reproduces seasonality but the normalised mean bias spans
\textbf{$-$68\% to $+$28\%} --- a factor-of-two-plus spread in a quantity inventories
report without error bars, attributed largely to unmonitored biomass burning in mainland
China. Residential burning across East Asia dominates; Southeast Asian transport matters
in summer and, via free-tropospheric pathways, \textbf{persists in winter} when
continental air masses are assumed to exclude it. \emph{Limitation}: three sites cannot
constrain a regional inventory, and $^{14}$C apportionment inherits uncertainty from the
assumed fossil/non-fossil end-members. Relevance: this is the cleanest current
demonstration that source-diagnostic \emph{measurement}, not finer model resolution, is
what bounds inventory error --- the argument for putting source-resolving capability into
distributed networks.}
{10.1021/acs.est.6c00449}

\paperentry{Deng et al.\ 2026}
{Impact of black carbon on daytime valley and slope winds in idealised simulations}
{Atmos Chem Phys}
{High-resolution idealised WRF-Chem experiments, run with and without radiatively
interactive black carbon, on thermally driven valley circulation. Absorption warms the
upper boundary layer and cools lower levels, increasing stability, reducing surface
heating, and leaving up-slope winds \textbf{weaker and confined to a shallower layer}.
The feedback is self-reinforcing: the circulation that ventilates a valley is suppressed
by the aerosol it would otherwise disperse. \emph{Limitation}: two idealised simulations
with prescribed BC, no observational constraint, no real terrain --- the sign of the
effect is credible, the magnitude is not transferable. Relevance to monitoring: in
mountain basins the relationship between emissions and measured concentration is not
linear, and siting a network to sample a circulation that the pollution itself modifies
requires the vertical dimension, not a denser surface grid.}
{10.5194/acp-26-11171-2026}

\paperentry{\textit{Environ Sci Atmos} 2026a}
{Multi-scale variability and lagged meteorological controls of PM$_{2.5}$ in Delhi}
{Environ Sci Atmos}
{Wavelet coherence and causality analysis of Delhi PM$_{2.5}$ against meteorological
drivers, decomposing the record by timescale rather than fitting a single regression.
\textbf{The abstract is truncated at source} --- it stops mid-sentence after stating the
aim, and no coherence bands, lag lengths, data period or station count are recoverable.
What can be said is that the method is the right one for the problem: Delhi's PM$_{2.5}$
is driven by processes on separable scales (diurnal boundary-layer collapse, synoptic
ventilation, seasonal crop-residue burning) that a time-domain model conflates, and
lagged causality is what forecasting a distributed network needs. Carried for the
register with its numbers unstated. Re-check on the weekly abstract sweep.}
{10.1039/d6ea00071a}

\paperentry{Li et al.\ 2026a}
{Upgrading diesel vehicle after-treatment technology significantly changed the
composition of reactive nitrogen in exhaust}
{Environ Sci Technol}
{Chassis-dynamometer and controlled testing of light-duty diesel vehicles across China
IV, V and VI standards. Nitrogen speciation shifts fundamentally: \textbf{NO exceeds 80\%
of reactive nitrogen under China IV}, while NH$_3$ and N$_2$O become prominent under
China VI. Mechanisms are attributed component-wise --- diesel oxidation catalysts raise
the NO$_2$/NO$_x$ ratio, SCR produces NH$_3$ slip, ammonia-slip catalysts generate N$_2$O.
SHAP analysis of tree models identifies emission standard, exhaust temperature and vehicle
speed as dominant drivers. \emph{Limitation}: no fleet sample size or vehicle count is
given, and dynamometer cycles under-represent cold starts and real-world transients where
SCR is least effective. Relevance: NH$_3$ is a secondary-PM precursor and N$_2$O a
greenhouse gas, so successive NO$_x$ controls have moved the problem rather than closed
it. A traffic-adjacent network measuring NO$_x$ alone now under-reports the emission it
is meant to characterise.}
{10.1021/acs.est.6c00994}

\paperentry{Lin et al.\ 2026}
{Three-millennia alpine-lake records of divergent atmospheric Cd, Pb, Sb and Hg pollution
in southwestern China}
{Environ Sci Technol}
{Radiometrically dated sediment cores from three alpine lakes at $\sim$4000\,m in Yunnan
reconstruct 3000 years of atmospheric Sb, Pb, Cd and Hg deposition. Onset is
metallurgy-specific and spatially staggered: Sb and Pb from the mid-1400s in northeastern
Yunnan, $\sim$200 years later in the west with Cu/Ag smelting; Cd only after the early
1700s with Zn and Cu production. Preindustrial peaks fall between the early 1800s and
1900\,CE. \textbf{Since the 1950s Cd has reached unprecedented levels while Sb and Pb
remain below their preindustrial maxima}, and Hg rises monotonically from the mid-1400s.
\emph{Limitation}: lake sediment integrates catchment erosion as well as atmospheric
deposition, and the separation depends on normalisation assumptions not visible in the
abstract. Relevance: a 3000-year baseline is the only way to know what ``background''
means for a metal, and metal-resolved PM is where low-cost sensing has no capability at
all.}
{10.1021/acs.est.5c17960}

\paperentry{\textit{Atmos Environ} 2026c \textsc{metadata only}}
{Source-specific black carbon responses to short-term interventions and cross-year
emission changes in coastal Qingdao}
{Atmos Environ}
{\textsc{metadata only} --- no abstract retrievable. Carried because source-specific BC
response to \emph{short-term} interventions is a natural-experiment design, and coastal
sites separate shipping from land-based combustion in a way inland sites cannot. Read
alongside Chen et al.\ in this issue, which finds inventory-based BC attribution biased by
$-$68\% to $+$28\%: an intervention study measures the response directly and sidesteps
that. Sample period, apportionment method and intervention definition are all unstated.}
{10.1016/j.atmosenv.2026.122299}

\paperentry{\textit{Atmos Pollut Res} 2026 \textsc{metadata only}}
{High-resolution CFD modelling of traffic-related air pollutant dispersion in an urban
area of Istanbul}
{Atmos Pollut Res}
{\textsc{metadata only} --- no abstract in Crossref, Europe PMC or Semantic Scholar.
Carried because building-resolving CFD is the reference against which cheap dispersion
proxies --- Gaussian plume, land-use regression, distance-to-road --- are judged, and
because street-canyon gradients at the scale CFD resolves are exactly what a dense
low-cost network claims to capture and what a regulatory monitor cannot. Whether the model
is validated against measurement, and at what resolution, is unknown. Listed, not
weighted.}
{10.1016/j.apr.2026.103169}

\band{Occupational \& indoor\hfill 3 records}{Clay}

\paperentry{Alqahtani et al.\ 2026}
{Biomarker-based assessment of tobacco smoke exposure in Saudi females via urinary
cotinine}
{PLoS One}
{LC-MS/MS quantification of urinary cotinine in 133 Saudi women recruited May--December
2024, validated to ICH-M10 (linearity 2--500\,ng/mL, $R^2\geq0.9992$, accuracy
90--107\%, CV $<$10\%), with a population-specific cut-off of \textbf{7\,ng/mL}. Cotinine
was detected in \textbf{48 of 133} samples although \textbf{88.7\% self-reported as
non-smokers}, and 45.9\% reported household second-hand smoke. \emph{Limitations}: $n=133$
is small for establishing a population cut-off; the detected range (0.02--7.693\,ng/mL)
sits almost entirely \emph{below} the proposed 7\,ng/mL threshold, so the cut-off
classifies nearly every detected sample as non-smoking and its derivation is not
defensible from the reported distribution; and cotinine's short half-life indexes recent
exposure only. Relevance: a clean demonstration that self-report and biomarker diverge
sharply under social desirability pressure --- the same argument for objective personal
measurement over questionnaires that motivates wearable PM sensing.}
{10.1371/journal.pone.0355166}

\paperentry{Herkert et al.\ 2026}
{Contribution of indoor air and dust to children's PFAS exposure in the Faroe Islands}
{Environ Sci Technol}
{Measurement of nonvolatile precursors, volatile neutral precursors and legacy ionic
PFAS in air and dust across \textbf{40 Faroese homes}, with estimated daily intake
computed for 15-year-olds. EDI ranges 0.001 / 0.012 / 0.129\,ng/kg\,bw across low,
intermediate and high scenarios, equal to roughly \textbf{0.2\%, 2\% and 20\% of the EFSA
tolerable weekly intake}, against $\sim$60\% from pilot-whale consumption. Indirect
exposure to perfluoroalkyl carboxylic acids via air dominates, followed by PFOS in dust.
\emph{Limitations}: scenario-based EDIs are not measured intakes, and the 40-home sample
cannot support the tail of the distribution where the 20\% figure sits. Relevance: the
indoor compartment supplies a fifth of a regulatory reference dose in the high scenario
for a population with no local source --- a direct case for indoor sampling as a routine
part of exposure assessment rather than a specialist add-on.}
{10.1021/acs.est.6c02003}

\paperentry{Iavazzo et al.\ 2026}
{Particle characterization in an urban synthetic turf: turf and vehicular wear}
{Environ Geochem Health}
{XRF, SEM-EDS and dynamic light scattering on dust, turf fibre and crumb-rubber infill
from an urban synthetic pitch, with stub-based sampling worn \emph{on players}. Irregular
and elongated fine, submicron and nano-sized particles were observed; Fe dominated by
XRF, with Fe, Si, Al and Br-bearing particles by SEM-EDS and \textbf{Ni, Ag and Rh}
recovered from the player-worn stubs --- Rh being a clear three-way-catalyst tracer, so
the pitch is receiving vehicular as well as turf-wear particles. Ambient PM$_{2.5}$
averaged \textbf{10.93\,$\mu$g/m$^3$} (24-h) and O$_3$ reached 41.12\,ppb.
\emph{Limitations}: one field, no personal inhaled dose, and the deterministic USEPA
assessment --- which found no significant non-carcinogenic metal risk --- is mass-based
and cannot speak to nanoparticle toxicity. Relevance: the contrast between an
unremarkable regulatory 24-h mean and a demonstrably loaded microenvironment is the
clearest argument in this issue for personal and microenvironmental sampling.}
{10.1007/s10653-026-03407-1}

\band{Respiratory \& allergic\hfill 3 records}{Sage}

\paperentry{Bogunia-Kubik et al.\ 2026}
{Influence of environmental pollution on basophil reactivity to birch pollen allergens}
{Allergy}
{Fifty birch-allergic patients and 25 non-atopic controls followed over the 2023 and 2024
pollen seasons in Krak\'ow, with app-recorded daily nasal, ocular and asthma symptoms
against continuously monitored PM$_{2.5}$, PM$_{10}$, O$_3$ and birch pollen. Basophil
activation tests used commercial extract, recombinant Bet\,v\,1, and \textbf{natural
pollen collected from highly and less polluted sites}. Pollen from polluted sites induced
significantly stronger CD63 expression ($p<0.0001$), and asthma symptoms correlated with
both pollen and PM ($p<0.05$, mainly 2023). \emph{Limitations}: the central comparison is
site-level, so pollution is confounded with everything else that differs between the two
collection sites; $n=75$ over two seasons with a season-specific correlation is thin; and
CD63 is an ex-vivo surrogate. Relevance: if pollution modifies allergen \emph{potency}
rather than only co-occurring with it, exposure models that treat pollen and PM as
independent additive terms are mis-specified.}
{10.1111/all.70474}

\paperentry{Govindaraju et al.\ 2026}
{Protective role of dietary antioxidant intake on long-term effects of extreme PM$_{2.5}$
exposure on respiratory health}
{J Occup Environ Med}
{Cross-sectional analysis of 448 Hazelwood Health Study adults, with individual PM$_{2.5}$
from the 2014 Morwell coal-mine fire reconstructed by combining modelled fire-related
concentrations with contemporaneous time--location diaries --- a better exposure
assignment than most disaster epidemiology achieves. Vitamins A and E, magnesium and zinc
attenuated the PM$_{2.5}$--chronic-cough association, and all but magnesium the
phlegm association; omega-3 attenuated the PM$_{2.5}$--COPD association. \emph{Limitations
are severe}: diet was measured cross-sectionally years after a 2014 exposure, so reverse
causation and healthy-diet confounding are unaddressed; \textbf{no effect sizes or
intervals are reported in the abstract}; and multiple testing across several nutrients
and several outcomes is not accounted for. Relevance: effect-modification analyses of this
kind demand exposure precision that only time-resolved personal or
time-activity-weighted measurement supplies.}
{10.1097/JOM.0000000000003856}

\paperentry{Wang et al.\ 2026a}
{COVID-19 social interventions amplified indoor-source pollutant--respiratory
associations in Xi'an}
{Front Public Health}
{Six-year time series of \textbf{2{,}052{,}777} respiratory outpatient visits over 2192
days at a tertiary hospital in Xi'an (2020--2025), stratified into strict control,
transition and post-pandemic phases, with quasi-Poisson GAMs and pollutant-by-phase
interaction terms. Full-period associations were strongest for SO$_2$ (lag 4, ER
\textbf{9.11\%}, 95\% CI 3.86--14.64) and NO$_2$ (lag 5, \textbf{2.72\%}, 2.05--3.40).
Phase-stratified, SO$_2$ was \textbf{18.40\% (10.29--27.11)} under strict control versus
\textbf{2.91\% ($-$4.82 to 11.28)} post-pandemic; PM$_{2.5}$ and O$_3$ showed no
interaction ($p>0.9$). \emph{The authors state plainly that the interaction tests do not
reach significance} ($p=0.13$, $p=0.15$), and single-hospital outpatient counts confound
exposure effects with healthcare-seeking behaviour that the pandemic also changed.
Relevance: the proposed indoor-amplification mechanism is unfalsifiable with ambient
monitors alone, which is the argument for indoor and personal measurement in
time-series designs.}
{10.3389/fpubh.2026.1896806}

\band{Neuro / mental health\hfill 2 records}{Deep}

\paperentry{Sullivan et al.\ 2026}
{Longitudinal impacts of childhood neighborhood air pollution on white matter integrity
and cognition}
{Dev Psychopathol}
{ABCD Study, $n=4{,}511$ children (47.4\% female, baseline mean age 9.9), with baseline
PM$_{2.5}$ tested against fronto-limbic fractional anisotropy and mean diffusivity across
follow-up, controlling for neighbourhood factors. \textbf{The direction of association
reverses with age}: at 9, higher PM$_{2.5}$ predicted \emph{increased} FA in anterior
thalamic radiation, fornix and uncinate; by 14, \emph{reduced} FA in all tracts. In
superior longitudinal fasciculus, higher MD at 9 and reduced FA at 12 and 14.
\textbf{Associations with cognition and substance use were null.} \emph{Limitations}:
baseline-only exposure ignores subsequent trajectory; neighbourhood-level assignment
carries substantial misclassification; and a sign reversal is equally consistent with
accelerated early myelination followed by decline, or with residual confounding that
changes with age. Relevance: a design that needs longitudinal, individually resolved
exposure and had only a single neighbourhood-level baseline value --- the precise gap
distributed personal sensing addresses.}
{10.1017/S095457942610159X}

\paperentry{\textit{Parkinsons Dis} review 2026}
{Air pollution and Parkinson's disease pathology: clinical evidence and molecular
mechanisms}
{Parkinsons Dis}
{Narrative review linking long-term NO$_x$, PM and O$_3$ exposure to Parkinson's incidence
and to worse prognosis, with proposed mechanisms spanning gut-microbiome alteration,
oxidative stress, neuroinflammation and promotion of $\alpha$-synuclein aggregation.
\emph{No pooled estimate, no systematic search strategy, no risk-of-bias assessment and
no quantitative synthesis} --- it is a mechanistic orientation piece, not evidence.
Tier C accordingly. The gut-microbiome route is the one worth noting, because it implies
an ingestion and mucosal pathway for deposited particles that inhaled-dose models do not
represent, and because it predicts that the relevant exposure metric is cumulative
deposited mass in a specific size range rather than ambient concentration. Carried for
the register and for that single mechanistic pointer.}
{10.1155/padi/6914112}

\band{Other clinical endpoints\hfill 2 records}{Slate}

\paperentry{\textit{Curr Hematol Malig Rep} 2026}
{Linking air pollution and hematologic malignancies: from epidemiology to molecular
mechanisms}
{Curr Hematol Malig Rep}
{Review of outdoor air pollution and haematologic malignancy across acute and chronic
leukaemias, lymphomas, MDS, CHIP and myeloma, implicating PM, benzene, NO$_2$, SO$_2$,
arsenic and ethylene oxide through oxidative stress, DNA damage, epigenetic dysregulation,
chronic inflammation and haematopoietic stem-cell dysfunction. Evidence is graded
strongest for AML with benzene; ALL associations with traffic-related pollution, NO$_2$
and PM are described as less consistent; evidence for chronic leukaemias and myeloma is
limited. \emph{Not systematic}: no search protocol, no pooled estimates, no bias
assessment, and the authors themselves flag exposure misclassification and residual
confounding. The CHIP link is the forward-looking one --- a clonal intermediate would give
this literature a measurable early endpoint. Relevance: benzene and PM co-vary in traffic
sources, and separating them requires speciated measurement that mass monitors cannot
provide.}
{10.1007/s11899-026-00785-2}

\paperentry{Wang et al.\ 2026b}
{Ambient PM$_{2.5}$ and thyroid health: 24{,}333 participants in Zhejiang, China}
{J Hazard Mater}
{Community-based cross-section of \textbf{24{,}333} participants (2022--2025) spanning
minors and adults, with PM$_{2.5}$ assigned from geocoded residential addresses and mixed
models estimating associations with FT4, TSH, subclinical hypothyroidism and autoimmune
thyroiditis. Per 10\,$\mu$g/m$^3$: FT4 \textbf{$-$3.57\% (95\% CI $-$4.39, $-$2.75)}, TSH
\textbf{$+$12.15\% (7.92, 16.56)}, subclinical hypothyroidism OR \textbf{1.27 (1.06,
1.51)}; no clear association with autoimmune thyroiditis. The FT4 effect was larger in
under-18s and in BMI $<$18.5; the TSH and hypothyroidism effects larger in those
$\geq$50 and BMI $\geq$24. \emph{Limitations}: cross-sectional, so exposure precedence is
assumed not observed; residential-address assignment ignores time--activity; and iodine
status is invoked as a rationale but its adjustment is not described in the abstract.
The internally coherent FT4-down/TSH-up pattern argues against pure chance. Relevance: the
paediatric and low-BMI susceptibility windows require exposure resolved at the individual
rather than the address.}
{10.1016/j.jhazmat.2026.143231}

\band{Mechanistic toxicology\hfill 1 record}{Violet}

\paperentry{Ichinose et al.\ 2026}
{Time-dependent BALT formation and pulmonary inflammation induced by different
particulate matter types}
{J Hazard Mater}
{Murine time-course comparing TiO$_2$ nanoparticles, diesel exhaust particles and Asian
sand dust, with and without ovalbumin. All three supported bronchus-associated lymphoid
tissue formation, more strongly with allergen. With allergen, an early IL-1$\alpha$ rise
preceded type-2 responses and BALT expansion; \textbf{without allergen, BALT still formed
via CXCL1 and neutrophil accumulation}, most clearly after repeated TiO$_2$ --- two
distinct routes to the same structure. 3D reconstruction showed tertiary lymphoid
structures extending from large airways distally, and TiO$_2$ was localised in alveoli,
alveolar macrophages, BALT and hilar lymph nodes. \emph{Limitations}: instilled or
high-dose murine exposure is not ambient inhalation; no dose--response to ambient
concentrations; and the three particle types differ in dose, size and surface chemistry
simultaneously. Relevance: particle \emph{identity}, not mass, determines which
inflammatory route is taken --- the standing mechanistic case for composition-resolved
monitoring.}
{10.1016/j.jhazmat.2026.143239}

\band{Burden, policy \& mitigation\hfill 3 records}{Coral}

\paperentry{Lee et al.\ 2026}
{Toward a unified framework for climate and health research: gaps in exposure metrics and
health outcomes}
{J Korean Med Sci}
{Mapping review of \textbf{611} original human studies (PubMed, Cochrane, Embase;
January 2000 -- June 2024) on climate-related exposures and health. The distribution is
sharply skewed: \textbf{526 studies on temperature and temperature-related hazards, 88 on
precipitation, and only 39 on air pollution}, with extreme weather barely represented.
Outcomes concentrate in all-cause mortality (249), cardiovascular disease (159) and
communicable disease (148), while neurodegenerative disease (8) and endocrine disorders
(4) are almost absent. Exposure-metric definitions are described as substantially
heterogeneous. \emph{Limitations}: counts are of studies, not of evidence quality, and
English-only inclusion biases the geography. Relevance: the two endpoints this issue
adds --- thyroid function and fronto-limbic white matter --- sit in the two thinnest
outcome categories the review identifies, and the metric heterogeneity it documents is the
same problem that makes cross-study pooling of PM exposure estimates unreliable.}
{10.3346/jkms.2026.41.e294}

\paperentry{Okeke et al.\ 2026}
{Agricultural methane emissions and non-communicable disease mortality in West Africa}
{Environ Monit Assess}
{Ecological panel of 11 West African countries, 2015--2020, relating agricultural methane
(Mt CO$_2$e, World Bank WDI) to the probability of dying from NCDs between ages 30 and 70,
using pooled OLS, country fixed effects and mixed models, with The Gambia
($-$28.3\% emissions) as a descriptive contrast. \emph{This is a weak design and should be
read as such}: 11 countries $\times$ 6 years is 66 country-years, the exposure is a
national accounting aggregate with no measured concentration and no plausible direct
mechanism to NCD mortality, and the outcome is a modelled WHO indicator rather than
observed deaths. Ecological fallacy and confounding by development trajectory are
unaddressed. Tier C, carried only because it is one of the few Sub-Saharan records in
this watch and marks where the climate--health literature is expanding without
measurement infrastructure --- which is itself the finding.}
{10.1007/s10661-026-15800-5}

\paperentry{\textit{Atmos Environ} 2026d \textsc{metadata only}}
{Drivers of air pollutant and greenhouse gas emissions since 2010 in the Greater Bay
Area, China}
{Atmos Environ}
{\textsc{metadata only} --- no abstract retrievable at build time. Carried for the
register because decomposition of co-emitted air pollutants and greenhouse gases over a
15-year window in a single megacity cluster is the analysis that determines whether
observed PM improvement came from control technology, fuel switching, activity shift or
economic structure --- and those four have different implications for what happens next.
Method, pollutant coverage and decomposition scheme are all unstated. Lowest-weight
record in the issue.}
{10.1016/j.atmosenv.2026.122301}

\clearpage

% ==================================================================== TABLE
\section*{\sffamily\bfseries\color{Deep}Corpus register}
\vspace{-1.5mm}{\color{Amber}\rule{\linewidth}{1.1pt}}\vspace{2.5mm}

{\fontsize{6.7}{8.4}\selectfont
\setlength{\tabcolsep}{2.4pt}
\begin{longtable}{@{}>{\RaggedRight}p{25mm}>{\RaggedRight}p{41mm}>{\RaggedRight}p{27mm}%
>{\RaggedRight}p{23mm}>{\RaggedRight}p{18mm}>{\RaggedRight}p{15mm}@{}}
\toprule
\textsf{\textbf{First author}} & \textsf{\textbf{Focus}} & \textsf{\textbf{Journal}} &
\textsf{\textbf{Design}} & \textsf{\textbf{Metric}} & \textsf{\textbf{Region}}\\
\midrule
\endfirsthead
\toprule
\textsf{\textbf{First author}} & \textsf{\textbf{Focus}} & \textsf{\textbf{Journal}} &
\textsf{\textbf{Design}} & \textsf{\textbf{Metric}} & \textsf{\textbf{Region}}\\
\midrule
\endhead
\multicolumn{6}{@{}l@{}}{\cellcolor{Mist}\textsf{\textbf{\textcolor{Deep}{Neuro / mental health}}}}\\[0.6mm]
\href{https://doi.org/10.1001/jamanetworkopen.2026.25015}{Dales} & Wildfire-sourced PM$_{2.5}$ and ED visits for migraine and other primary headache & \textit{JAMA Netw Open} & Case-crossover & PM$_{2.5}$ (wildfire) & Canada\\
\href{https://doi.org/10.1093/ageing/afag215}{Glover} & Umbrella meta-review of urban environmental exposures and cognitive health & \textit{Age Ageing} & Umbrella review / meta-analysis & PM$_{2.5}$ & Global\\
\href{https://doi.org/10.1016/j.nsa.2026.107019}{Goossens} & In utero air pollution exposure and pre-/postnatal brain development & \textit{Neurosci Appl} & Narrative review & PM$_{2.5}$ / PM$_{10}$ / PM$_{\mathrm{coarse}}$ & Global\\
\href{https://doi.org/10.1016/j.isci.2026.116617}{Ma} & PM$_{2.5}$-induced neurotoxicity and depression: exposure risks and mechanisms & \textit{iScience} & Narrative review & PM$_{2.5}$ + components & Global\\
\href{https://doi.org/10.3760/cma.j.cn112150-20250804-00752}{Pan} & Maternal PM exposure and offspring neurodevelopment: effects and mechanisms & \textit{Zhonghua Yu Fang Yi Xue Za Zhi} & Narrative review & PM (bulk) & China\\
\addlinespace[0.9mm]
\multicolumn{6}{@{}l@{}}{\cellcolor{Mist}\textsf{\textbf{\textcolor{Teal}{Cardiovascular \& metabolic}}}}\\[0.6mm]
\href{https://doi.org/10.3760/cma.j.cn112338-20251120-00843}{Huang} & PM$_{2.5}$ and dyslipidaemia in rural Chinese children; dietary effect modification & \textit{Zhonghua Liu Xing Bing Xue Za Zhi} & Cross-sectional (multistage) & PM$_{2.5}$ & China\\
\href{https://doi.org/10.32394/pe/218966}{Kazmierska} & Air pollution and cardiovascular disease: review of current data & \textit{Przegl Epidemiol} & Narrative review & PM$_{2.5}$ / PM$_{10}$ & Global\\
\href{https://doi.org/10.1007/s00125-026-06814-2}{Xi} & Air pollution, built environment and temperature vs. prediabetes phenotypes & \textit{Diabetologia} & Prospective cohort & PM$_{2.5}$ & Germany\\
\href{https://doi.org/10.1161/JAHA.125.046917}{Yang} & Short-term PM$_{2.5}$ and hospitalisation/ED use in atrial fibrillation & \textit{J Am Heart Assoc} & Case-crossover & PM$_{2.5}$ (wildfire-influenced) & USA\\
\addlinespace[0.9mm]
\multicolumn{6}{@{}l@{}}{\cellcolor{Mist}\textsf{\textbf{\textcolor{Sky}{Reproductive \& developmental}}}}\\[0.6mm]
\href{https://doi.org/10.1016/j.jes.2025.10.043}{Chen} & Prenatal air pollution, fetal growth trajectory phenotypes and child neurodevelopment & \textit{J Environ Sci (China)} & Prospective cohort & PM$_{2.5}$ / PM$_{10}$ & China\\
\href{https://doi.org/10.1016/j.ecoenv.2026.120566}{Huang} & Joint prenatal bisphenol + air-pollutant mixture exposure and preschool IQ & \textit{Ecotoxicol Environ Saf} & Prospective cohort (mixtures) & PM$_{2.5}$ / PM$_{10}$ & China\\
\href{https://doi.org/10.1007/s40572-026-00552-8}{Peterson} & Ambient air pollution, PAHs and pubertal development: critical review & \textit{Curr Environ Health Rep} & Critical review & PM$_{2.5}$ components & Global\\
\href{https://doi.org/10.1016/j.envres.2026.125313}{Xie} & Long-term PM$_{2.5}$ and timing of pubertal development in boys and girls & \textit{Environ Res} & Prospective cohort & PM$_{2.5}$ & China\\
\addlinespace[0.9mm]
\multicolumn{6}{@{}l@{}}{\cellcolor{Mist}\textsf{\textbf{\textcolor{Sage}{Respiratory \& allergic}}}}\\[0.6mm]
\href{https://doi.org/10.1007/s00484-026-03289-z}{Bayar Muluk} & Climate change and otolaryngology: airway disease impacts and emission reduction & \textit{Int J Biometeorol} & Narrative review & PM$_{2.5}$ & Global\\
\href{https://doi.org/10.4103/mgr.MEDGASRES-D-25-00164}{Liao} & Short-term air pollution, respiratory symptoms and lung function in asthma & \textit{Med Gas Res} & Panel study & PM$_{2.5}$ / PM$_{10}$ & China\\
\href{https://doi.org/10.3390/epidemiologia7040097}{Sanchez de Prada} & Ambient PM and asthma-coded admissions in children 0-3 y, Spain & \textit{Epidemiologia} & Ecological panel & PM$_{2.5}$ / PM$_{10}$ & Spain\\
\href{https://doi.org/10.1016/j.jhazmat.2026.143076}{Yang} & Allergenic vs. toxic risk of airborne seed hairs compared with ambient PM$_{2.5}$ & \textit{J Hazard Mater} & Field measurement & PM$_{2.5}$ + bioaerosol & China\\
\addlinespace[0.9mm]
\multicolumn{6}{@{}l@{}}{\cellcolor{Mist}\textsf{\textbf{\textcolor{Violet}{Mechanistic toxicology}}}}\\[0.6mm]
\href{https://doi.org/10.1016/j.jes.2025.12.074}{Ding} & PM$_{2.5}$ promotes vascular calcification via miR-27b-5p/HES1 signalling & \textit{J Environ Sci (China)} & In vitro + organoid + murine & PM$_{2.5}$ & China\\
\href{https://doi.org/10.1007/s00360-026-01698-5}{Maraschi} & Osmoregulatory disruption and metal bioaccumulation from settleable atmospheric PM & \textit{J Comp Physiol B} & Controlled exposure (ecotox) & Settleable APM & Brazil\\
\href{https://doi.org/10.1016/j.jes.2025.11.018}{Wu} & Macrophage NLRP3 inflammasome in real-ambient PM-induced hepatic glucose dysregulation & \textit{J Environ Sci (China)} & Real-ambient murine + co-culture & Real-ambient PM & China\\
\addlinespace[0.9mm]
\multicolumn{6}{@{}l@{}}{\cellcolor{Mist}\textsf{\textbf{\textcolor{Clay}{Occupational \& indoor}}}}\\[0.6mm]
\href{https://doi.org/10.3389/fmedt.2026.1796828}{Chen} & Nested suction instrument to cut intraoperative surgical-smoke PM$_{2.5}$ & \textit{Front Med Technol} & Bench / device experiment & PM$_{2.5}$ (surgical smoke) & China\\
\href{https://doi.org/10.1093/annweh/wxag055}{Cole} & Determinants of respirable crystalline silica exposure in Australian tunnelling & \textit{Ann Work Expo Health} & Observational exposure study & Respirable dust / RCS & Australia\\
\href{https://doi.org/10.1016/j.jhazmat.2026.143081}{Ding} & Stage-resolved hair Zn-isotope dynamics as a firefighting combustion-exposure tracer & \textit{J Hazard Mater} & Repeated-measures biomarker & Combustion-derived PM & China\\
\href{https://doi.org/10.3390/jox16040119}{Ramsden} & Xenobiotic hazards in aircraft cabin air & \textit{J Xenobiot} & Review + risk economics & Ultrafine particles & Global\\
\addlinespace[0.9mm]
\multicolumn{6}{@{}l@{}}{\cellcolor{Mist}\textsf{\textbf{\textcolor{Amber}{Sensing, forecasting \& instrumentation}}}}\\[0.6mm]
\href{https://doi.org/10.1016/j.jenvman.2026.130515}{Shi} & Adaptive spatiotemporal graph transformer for multi-site multi-step PM$_{2.5}$ forecasting & \textit{J Environ Manage} & Deep-learning model & PM$_{2.5}$ & China / India\\
\href{https://doi.org/10.3390/bios16070380}{Tsuchiya} & Seasonal and indoor-outdoor characteristics of fluorescent aerosol particles in offices & \textit{Biosensors} & Real-time sensor campaign & Bioaerosol / coarse & Japan\\
\addlinespace[0.9mm]
\multicolumn{6}{@{}l@{}}{\cellcolor{Mist}\textsf{\textbf{\textcolor{Amber}{Exposure assessment \& modelling}}}}\\[0.6mm]
\href{https://doi.org/10.1029/2025GH001621}{Solorzano-Araque} & Multiscale exposure modelling: LOTOS-EUROS + mobility surveys + particle toxicity & \textit{GeoHealth} & CTM + mobility integration & PM$_{2.5}$ / PM$_{10}$ & Colombia\\
\addlinespace[0.9mm]
\multicolumn{6}{@{}l@{}}{\cellcolor{Mist}\textsf{\textbf{\textcolor{Coral}{Burden, policy \& mitigation}}}}\\[0.6mm]
\href{https://doi.org/10.1002/hsr2.72886}{Karamzad} & Ischaemic heart disease burden in MENA 1990-2021; PM attributable fraction 28.2\% & \textit{Health Sci Rep} & GBD secondary analysis & Ambient PM & MENA\\
\href{https://doi.org/10.1371/journal.pone.0347989}{Rees} & HEPA filtration to reduce respiratory infection in care homes & \textit{PLoS One} & RCT process evaluation & Filtered PM & UK\\
\href{https://doi.org/10.1016/j.jes.2026.01.009}{Shan} & Policy-driven emission reductions and post-pandemic rebound, high-resolution inventories & \textit{J Environ Sci (China)} & Emission inventory & PM emissions & China\\
\addlinespace[0.9mm]
\multicolumn{6}{@{}l@{}}{\cellcolor{Mist}\textsf{\textbf{\textcolor{Slate}{Other clinical endpoints}}}}\\[0.6mm]
\href{https://doi.org/10.3390/dj14070440}{Lupita} & PM$_{2.5}$, preventive-care capacity and caries/edentulism burden across the EU & \textit{Dent J} & Ecological panel & PM$_{2.5}$ & EU-27\\
\href{https://doi.org/10.1016/j.jdin.2026.06.008}{Park} & Long-term PM exposure and risk of androgenetic alopecia & \textit{JAAD Int} & Nationwide cohort & PM$_{2.5}$ / PM$_{10}$ & South Korea\\
\href{https://doi.org/10.3171/2026.3.PEDS2615}{Sun} & Environmental exposure and area deprivation vs. paediatric hydrocephalus outcomes & \textit{J Neurosurg Pediatr} & Registry cohort + geospatial & PM$_{2.5}$ (satellite-derived) & USA\\
\addlinespace[0.9mm]
\bottomrule
\end{longtable}}


\vspace{4mm}
\begin{keybox}
{\sffamily\bfseries\footnotesize\color{Deep}Provenance and method}\\[1.2mm]
{\fontsize{7.2}{8.8}\selectfont
Entry window \textbf{2026-08-11 to 2026-08-11}, contiguous with the 10 August issue; no
gap. \textbf{PubMed} connector, three independent queries against \texttt{[Date - Entry]}:
health axis 9, instrumentation axis 10, broad MeSH sweep 23 --- \textbf{33 distinct}
records, \textbf{18 carried}. Ten of the 33 carry an Entrez create date of 10 August but
an \texttt{[EDAT]} of 11 August; none was in \texttt{seen.json}, so they belong to this
issue and the two date fields are not interchangeable. Local harvester: PubMed health 18,
PubMed sensing 1, Europe PMC \texttt{CREATION\_DATE} 1 (\textbf{1 carried}, a Research
Square preprint), Crossref by ISSN across 17 journals 12 (\textbf{7 carried}).
\textbf{Consensus} returned 20, top 3 retrievable on the free tier, \textbf{0 carried} ---
all three resolved to Crossref creation dates of 26 March 2026, 11 November 2025 and
21 April 2026 and all three were already in \texttt{seen.json} on DOI. That is the
\textbf{fifth consecutive run} confirming Consensus as a relevance index that cannot
resolve a single day; it continues to earn its place as a duplicate check, not as a
recency leg. \textbf{arXiv} 60 records, \textbf{1} published on or after 5 August and
off-topic (a climate--health bibliometric survey). \textbf{OpenAlex} HTTP 429,
\textbf{ninth consecutive run}. \textbf{ClinicalTrials.gov} returned \textbf{0} trials
registered or updated on 11 August; \texttt{trials.json} unchanged, and the trial watch is
surfaced in the weekly rollup, not here. \textbf{24 records rejected}, each with a reason
in \texttt{state/rejected.jsonl} --- the large majority environmental-science records
surfaced by the broad MeSH sweep with no atmospheric particulate pathway (freshwater
mercury, soil organic carbon, riverine methane, waterborne microplastics).
\textbf{Six carried records have no abstract} in Crossref, Europe PMC or Semantic
Scholar --- four Elsevier, one Springer, one RSC --- and are marked \textsc{metadata
only}; a seventh, the RSC Delhi wavelet paper, has an abstract truncated mid-sentence at
source. Fourth consecutive issue affected; the re-check remains a weekly sweep.
\textbf{Three defects were closed in this issue.} First and most serious:
\texttt{plots.py} carried a \textbf{silent hardcoded fallback for f6}, so an issue whose
corpus JSON lacked a \texttt{LIFECOURSE} key rendered life-course counts
\texttt{[5, 3, 2, 3, 4]} and exemplars (``firefighter biomarkers'', ``care-home
filtration'', ``PAH endocrine axis'') drawn from issues weeks earlier. Today's corpus
contains none of those records, so the figure would have shipped numbers describing no
paper in the issue. It was caught only by reading the rendered PNG against the store.
\texttt{corpus.save()} now takes and writes \texttt{lifecourse}, and a missing key is a
hard build failure rather than a plausible-looking placeholder. Second: the PubMed
\texttt{efetch} parser was reading DOIs with an unscoped \texttt{.//ArticleId} XPath,
which matched reference-list identifiers --- \textbf{15 of 33 records} received a DOI
belonging to a cited paper, several from the 2010s. The lookup is now scoped to
\texttt{PubmedData/ArticleIdList} with an \texttt{ELocationID} fallback, and
\texttt{check\_dois.py} passes 0 fail / 0 warn. Third: the geography mapper fell through
for \emph{Faroe Islands}, \emph{West Africa} and \emph{Idealised terrain}, silently
pooling all three into ``Global / multi-region'' and inflating that bar from 8 to 10; all
three are now mapped. \texttt{design\_group} did \textbf{not} fall through this issue ---
the first clean run in four --- because every design label written on a record was chosen
from the existing mapped vocabulary. The combined \texttt{f2\_design\_endpoint.png} panel
is \textbf{not used in this issue}: with 14 distinct endpoint labels its y-axis text
overlaps the donut centre and the architecture legend, so \texttt{f2a} and \texttt{f2b}
are shown separately pending a layout fix.
Subtopic, design, particle-metric and geography labels are assigned from title, abstract,
keywords and MeSH; assignment is single-label by dominant contribution and is a judgement,
not a controlled vocabulary. Effect estimates on the forest plot are heterogeneous by
metric and are not pooled: percentage excess risks are shown as $1+\mathrm{ER}/100$ and
percentage changes in continuous biomarkers as geometric-mean ratios, purely so they share
a ratio axis.
All DOI links resolve to the publisher of record. Content derived from PubMed, Crossref
and Europe PMC metadata; abstracts remain the copyright of their respective publishers.\par}
\end{keybox}

\end{document}
